% GENERATED FILE. Edit canonical lessons, then run npm run generate:books.
% canonical-source-hash: fnv1a64:bdde79dd15af58a7
% canonical-lessons: PT-C416-painel, PT-C416-atraso, PT-C416-atrasar, PT-C416-suprimir, PT-C416-transbordo, PT-R416-revisao-409-416

\chapter{Atrasos e transbordos}
\label{ch:pt-atrasos-e-transbordos}
\hlchaptermodality{416}

\begin{chapteropening}
\textbf{By the end of this chapter:} \emph{I can deal with delays and changes: the departures board, the delay, running late, a cancelled service and changing trains.}

Complete the last lesson of chapter 416: I can deal with delays and changes: the departures board, the delay, running late, a cancelled service and changing trains.
\end{chapteropening}

\section[o painel]{o painel — the board, the panel}

\label{lesson:PT-C416-painel}



\begin{quote}
Before the new one: say the Portuguese for the train driver, then the Portuguese for the engine.
\end{quote}

\subsection*{You'll want to know: o painel}
\textbf{o painel} — \textquotedblleft{}the board, the panel\textquotedblright{}.

\begin{quote}
O painel da estação mostra as partidas.
\end{quote}

\begin{quote}
\emph{The station board shows the departures.}
\end{quote}

\emph{O painel solar} — the solar panel. \emph{O painel} of a car is the dashboard.

\begin{cousinweb}
\emph{painel} came from Old French \emph{panel}, from Vulgar Latin \emph{pannellus}, a small form of Latin \emph{pannus} 'piece of cloth' (EtymDB).
\end{cousinweb}

\begin{grammarlens}[title={What you've built}]
One more word for telling, explaining and describing at B1.
\end{grammarlens}

\subsection*{Your turn}
\begin{itemize}
\raggedright
  \item \emph{Say it:} \emph{o painel}
  \item \emph{Say it:} \emph{o painel}, once more
  \item \emph{Say it:} a sentence of your own with \emph{o painel}
  \item \emph{Recall:} say the Portuguese for the train driver, then the Portuguese for the engine, then say \emph{o painel} again
\end{itemize}

\begin{tcolorbox}[breakable,colback=teal!4,colframe=teal!35!black,title={Before you move on}]
What does o painel mean? (\textquotedblleft{}The board, the panel\textquotedblright{}.) Say it once more.
\end{tcolorbox}
\section[o atraso]{o atraso — the delay}

\label{lesson:PT-C416-atraso}



\begin{quote}
Before the new one: say the Portuguese for the engine, then the Portuguese for the board.
\end{quote}

\subsection*{You'll want to know: o atraso}
\textbf{o atraso} — \textquotedblleft{}the delay\textquotedblright{}.

\begin{quote}
O comboio tem um atraso de vinte minutos.
\end{quote}

\begin{quote}
\emph{The train is twenty minutes late.}
\end{quote}

\emph{Estar atrasado} — to be late. \emph{Com atraso} — late.

\begin{cousinweb}
\emph{atraso} is built on \emph{atrasar}, which is built on \emph{atrás} 'behind' (from \emph{a} and \emph{trás}) with \emph{-ar} (EtymDB); \emph{trás} goes back to Latin \emph{trāns} 'across, beyond'.
\end{cousinweb}

\begin{grammarlens}[title={What you've built}]
One more word for telling, explaining and describing at B1.
\end{grammarlens}

\subsection*{Your turn}
\begin{itemize}
\raggedright
  \item \emph{Say it:} \emph{o atraso}
  \item \emph{Say it:} \emph{o atraso}, once more
  \item \emph{Say it:} a sentence of your own with \emph{o atraso}
  \item \emph{Recall:} say the Portuguese for the engine, then the Portuguese for the board, then say \emph{o atraso} again
\end{itemize}

\begin{tcolorbox}[breakable,colback=teal!4,colframe=teal!35!black,title={Before you move on}]
What does o atraso mean? (\textquotedblleft{}The delay\textquotedblright{}.) Say it once more.
\end{tcolorbox}
\section[atrasar]{atrasar — to delay, to make late; to put back (a clock)}

\label{lesson:PT-C416-atrasar}



\begin{quote}
Before the new one: say the Portuguese for the board, then the Portuguese for the delay.
\end{quote}

\subsection*{You'll want to know: atrasar}
\textbf{atrasar} — \textquotedblleft{}to delay, to make late; to put back (a clock)\textquotedblright{}.

\begin{quote}
A neve atrasou todos os voos.
\end{quote}

\begin{quote}
\emph{The snow delayed all the flights.}
\end{quote}

\emph{Atrasar-se} — to be late. In autumn people \emph{atrasam o relógio}.

\begin{cousinweb}
\emph{atrasar} is built on \emph{atrás} 'behind' (from \emph{a} and \emph{trás}) with \emph{-ar} (EtymDB); \emph{trás} goes back to Latin \emph{trāns} 'across, beyond'.
\end{cousinweb}

\begin{grammarlens}[title={What you've built}]
One more word for telling, explaining and describing at B1.
\end{grammarlens}

\subsection*{Your turn}
\begin{itemize}
\raggedright
  \item \emph{Say it:} \emph{atrasar}
  \item \emph{Say it:} \emph{atrasar}, once more
  \item \emph{Say it:} a sentence of your own with \emph{atrasar}
  \item \emph{Recall:} say the Portuguese for the board, then the Portuguese for the delay, then say \emph{atrasar} again
\end{itemize}

\begin{tcolorbox}[breakable,colback=teal!4,colframe=teal!35!black,title={Before you move on}]
What does atrasar mean? (\textquotedblleft{}To delay, to make late; to put back (a clock)\textquotedblright{}.) Say it once more.
\end{tcolorbox}
\section[suprimir]{suprimir — to cancel (a service); to suppress, to remove}

\label{lesson:PT-C416-suprimir}



\begin{quote}
Before the new one: say the Portuguese for the delay, then the Portuguese for to delay.
\end{quote}

\subsection*{You'll want to know: suprimir}
\textbf{suprimir} — \textquotedblleft{}to cancel (a service); to suppress, to remove\textquotedblright{}.

\begin{quote}
Por causa da greve, suprimiram três comboios.
\end{quote}

\begin{quote}
\emph{Because of the strike, they cancelled three trains.}
\end{quote}

On the board: \emph{comboio suprimido} — train cancelled.

\begin{cousinweb}
\emph{suprimir} is a learned borrowing from Latin \emph{supprimere} 'to press down, to hold back' (EtymDB), built on \emph{premere} 'to press'.
\end{cousinweb}

\begin{grammarlens}[title={What you've built}]
One more word for telling, explaining and describing at B1.
\end{grammarlens}

\subsection*{Your turn}
\begin{itemize}
\raggedright
  \item \emph{Say it:} \emph{suprimir}
  \item \emph{Say it:} \emph{suprimir}, once more
  \item \emph{Say it:} a sentence of your own with \emph{suprimir}
  \item \emph{Recall:} say the Portuguese for the delay, then the Portuguese for to delay, then say \emph{suprimir} again
\end{itemize}

\begin{tcolorbox}[breakable,colback=teal!4,colframe=teal!35!black,title={Before you move on}]
What does suprimir mean? (\textquotedblleft{}To cancel (a service); to suppress, to remove\textquotedblright{}.) Say it once more.
\end{tcolorbox}
\section[o transbordo]{o transbordo — the change (of trains or buses), the transfer}

\label{lesson:PT-C416-transbordo}



\begin{quote}
Before the new one: say the Portuguese for to delay, then the Portuguese for to cancel.
\end{quote}

\subsection*{You'll want to know: o transbordo}
\textbf{o transbordo} — \textquotedblleft{}the change (of trains or buses), the transfer\textquotedblright{}.

\begin{quote}
Esta viagem tem um transbordo a meio do caminho.
\end{quote}

\begin{quote}
\emph{This trip has a change of trains halfway.}
\end{quote}

\emph{Fazer transbordo} — to change trains.

\begin{cousinweb}
\emph{transbordo} is built from \emph{trans-} 'across' and \emph{bordo} 'ship's side': first moving cargo from one ship to another. \emph{Bordo} goes back to Old Portuguese \emph{bordo} (EtymDB) and is usually traced to a Germanic word for 'plank, side'.
\end{cousinweb}

\begin{grammarlens}[title={What you've built}]
One more word for telling, explaining and describing at B1.
\end{grammarlens}

\subsection*{Your turn}
\begin{itemize}
\raggedright
  \item \emph{Say it:} \emph{o transbordo}
  \item \emph{Say it:} \emph{o transbordo}, once more
  \item \emph{Say it:} a sentence of your own with \emph{o transbordo}
  \item \emph{Recall:} say the Portuguese for to delay, then the Portuguese for to cancel, then say \emph{o transbordo} again
\end{itemize}

\begin{tcolorbox}[breakable,colback=teal!4,colframe=teal!35!black,title={Before you move on}]
What does o transbordo mean? (\textquotedblleft{}The change (of trains or buses), the transfer\textquotedblright{}.) Say it once more.
\end{tcolorbox}
\section[(dialogue)]{Revisão dos capítulos 409–416}

\label{lesson:PT-R416-revisao-409-416}



\begin{quote}
Say the Portuguese for the traffic jam and to slow down.
\end{quote}

\subsection*{Your turn}
Say these aloud:

\begin{itemize}
\raggedright
  \item the traffic jam, to slow down, the collision, to crash into, the diversion
  \item the bend, the speed bump, the shortcut, tarmac, to wind
  \item the hire, the driver, the number plate, the mileage, the excess
  \item the moped, the scooter, the saddle, to pedal, the parking meter
  \item the shelter, the terminus, the coach, the route, to get off
  \item the ticket office, the travel pass, to validate, the ticket inspector, the fine
  \item the carriage, the compartment, the bunk, the train driver, the engine
  \item the board, the delay, to delay, to cancel, the change
\end{itemize}

\begin{tcolorbox}[breakable,colback=teal!4,colframe=teal!35!black,title={Before you move on}]
The traffic jam? (\textbf{o engarrafamento}.) To slow down? (\textbf{abrandar}.) The collision? (\textbf{a colisão}.) To crash into? (\textbf{embater}.) The diversion? (\textbf{o desvio}.) The bend? (\textbf{a curva}.) The speed bump? (\textbf{a lomba}.) The shortcut? (\textbf{o atalho}.) Tarmac? (\textbf{o alcatrão}.) To wind? (\textbf{serpentear}.) The hire? (\textbf{o aluguer}.) The driver? (\textbf{o condutor}.) The number plate? (\textbf{a matrícula}.) The mileage? (\textbf{a quilometragem}.) The excess? (\textbf{a franquia}.) The moped? (\textbf{a motorizada}.) The scooter? (\textbf{a trotinete}.) The saddle? (\textbf{o selim}.) To pedal? (\textbf{pedalar}.) The parking meter? (\textbf{o parquímetro}.) The shelter? (\textbf{o abrigo}.) The terminus? (\textbf{o terminal}.) The coach? (\textbf{a camioneta}.) The route? (\textbf{o itinerário}.) To get off? (\textbf{apear}.) The ticket office? (\textbf{a bilheteira}.) The travel pass? (\textbf{o passe}.) To validate? (\textbf{validar}.) The ticket inspector? (\textbf{o revisor}.) The fine? (\textbf{a coima}.) The carriage? (\textbf{a carruagem}.) The compartment? (\textbf{o compartimento}.) The bunk? (\textbf{o beliche}.) The train driver? (\textbf{o maquinista}.) The engine? (\textbf{a locomotiva}.) The board? (\textbf{o painel}.) The delay? (\textbf{o atraso}.) To delay? (\textbf{atrasar}.) To cancel? (\textbf{suprimir}.) The change? (\textbf{o transbordo}.)
\end{tcolorbox}
